% !PRINTED 22
\weberpage{22}
With the help of the covariants we have reduced the form $f(x)$ to the
\gesperrt{canonical form}
\[
\varphi(y) = y_1^3+y_2^3+y_3^3+6my_1y_2y_3
\]

We have, if $r$ denotes the substitution determinant, set:
\begin{align*}
P &= \frac{(2+m^3)m^3f^2+(2-5m^3)mr^2f\Delta+3m^2r^4\Delta^4-r^6J}{(1+8m^3)^2}, % [sic] printed \Delta^4; the later specialization uses \Delta^2
\\
Q &= \frac{(1+2m^3)f-6mr^2\Delta}{1+8m^3}, \\
R &= \frac{m^2f+r^2\Delta}{1+8m^3}
\end{align*}
and obtained $y_1^3$, $y_2^3$, $y_3^3$ as roots of the cubic
equation
\[
u^3-Qu^2+Pu-R^3=0,
\]
and the discriminant of this cubic equation is
\[
D_1 = \frac{r^{18}K^2}{(1+8m^3)^6}.
\]

If one expresses this discriminant in terms of $P$, $Q$, $R$, one
obtains $K^2$ represented rationally in terms of $f$, $\Delta$, $J$. At
the passage of the Algebra mentioned we did not give this expression
explicitly. Now, however, we must form it, even if only under the
assumption $f=0$, i.\,e. only for the points of the curve.
But we have [according to Bd.\,I, § 50, (10)]
\[
D_1 = P^2Q^2+18PQR^3-4P^3-4Q^3R^3-27R^6.
\]
In this one has to set:
\[
P = \frac{3m^2r^4\Delta^2-r^6J}{(1+8m^3)^2}, \quad Q = \frac{-6mr^2\Delta}{1+8m^3}, \quad R = \frac{r^2\Delta}{1+8m^3},
\]
and one obtains, by a calculation that is not difficult, if one denotes
by $S$ and $T$ the two \gesperrt{invariants} of the curve of the third order
(Bd.\,II, § 108),
\[
(1+8m^3)^6D_1 = r^{18}(4J^3+108SJ\Delta^4-27T\Delta^6),
\]
and consequently
\weq{8}{K^2 = 4J^3+108SJ\Delta^4-27T\Delta^6.}
This equation, however, is not satisfied identically, but only under
the assumption $f=0$. The complete expression of $K^2$ in terms of
$J$, $\Delta$, $f$ becomes very much more complicated.
